\chapter{फूरिये श्रेणियाँ}\label{ch:b2:fourier}

क्या हर आवर्ती संकेत को शुद्ध ज्याओं और कोज्याओं से फिर से खड़ा किया जा
सकता है? फूरिये के साहसिक ``हाँ'' ने विश्लेषण की एक पूरी सदी रच दी। यह
अध्याय इस स्तर पर पहुँच में आने वाले दोनों स्तंभ सिद्ध करता है:
\emph{डिरिक्ले की प्रमेय} (\omterm{lem:b2:fourier:kernel}{डिरिक्ले कर्नेल} के द्वारा, खंडशः $C^1$ फलनों के
लिए \omterm{def:b2:funcseq:def}{बिंदुवार} पुनर्निर्माण) और \emph{पारसेवाल की सर्वसमिका} (किसी संकेत की
ऊर्जा उसकी संनादियों की ऊर्जाओं का योग होती है); और फिर क्लासिक संख्यात्मक
श्रेणियाँ काटता है --- सबसे पहले बाज़ल का $\sum 1/n^2 =
\pi^2/6$।

आगे सर्वत्र फलन $2\pi$-आवर्ती, खंडशः \omterm{def:b2:metric:continuity}{संतत} और सम्मिश्र-मान वाले हैं; और $\mathcal{C}$
\omterm{def:b2:metric:continuity}{संतत} फलनों को दर्शाता है।

\section{फूरिये गुणांक}

\begin{definition}\label{def:b2:fourier:coefficients}
$f$ के \emph{फूरिये गुणांक}\index{फूरिये गुणांक}
\[
c_n(f) = \frac{1}{2\pi}\int_{-\pi}^{\pi} f(t)\,\eu^{-\iu n t}\,\dd
t \qquad (n \in \Z),
\]
हैं, तथा वास्तविक-रूप गुणांक $a_n = c_n + c_{-n}$, $b_n = \iu(c_n
- c_{-n})$, जिससे \emph{फूरिये आंशिक
योग}\index{फूरिये आंशिक योग}
\[
S_N(f)(t) = \sum_{n=-N}^{N} c_n(f)\,\eu^{\iu nt}
= \frac{a_0}{2} + \sum_{n=1}^{N}\bigl(a_n\cos nt + b_n \sin
nt\bigr).
\]
बनते हैं। $\mathcal{C}$ पर \omterm{def:b2:hermitian:def}{हर्मीशियन अंतर्गुणन} $\langle f,
g\rangle = \frac{1}{2\pi}\int_{-\pi}^{\pi}\conj f\,g$ परिभाषित कीजिए: तब चरघातांकी $e_n(t) = \eu^{\iu nt}$
\emph{प्रसामान्य लांबिक} हैं ($\langle e_m, e_n\rangle = \delta_{mn}$, सीधा परिकलन), और $c_n(f) = \langle e_n, f\rangle$: अर्थात् फूरिये
विश्लेषण अनंत विमा में \omterm{def:b2:hermitian:adjoint}{हर्मीशियन} ज्यामिति (\cref{ch:b2:hermitian}) ही है।
\end{definition}

\begin{proposition}[बेसेल असमिका]\label{prop:b2:fourier:bessel}
$S_N(f)$ $f$ का घात $\leq N$ वाले त्रिकोणमितीय बहुपदों की समष्टि $\mathcal{T}_N$ पर लांबिक
प्रक्षेप है, और\index{बेसेल असमिका}
\[
\sum_{n=-N}^{N} \abs{c_n(f)}^2 \leq \norm f_2^2 =
\frac{1}{2\pi}\int_{-\pi}^{\pi} \abs f^2 :
\]
अतः श्रेणी $\sum \abs{c_n}^2$ अभिसरित होती है, और $\abs n \to \infty$ होने पर $c_n(f) \to 0$ (गुणांकों के लिए
रीमान--लेबेग)।
\end{proposition}

\begin{proof}
$f - S_N(f)$ हर $e_k$, $\abs k \leq N$ के लांबिक है ($\langle
e_k, f - S_N f\rangle = c_k - c_k = 0$): अतः $S_Nf$ $\mathcal{T}_N = \operatorname{Vect}(e_{-N}, \dots,
e_N)$ पर लांबिक प्रक्षेप है
(प्रथम वर्ष के खंड की प्रक्षेप प्रमेय, \omterm{def:b2:hermitian:adjoint}{हर्मीशियन} परिवेश में अक्षरशः)।
पाइथागोरस: $\norm f_2^2 = \norm{S_Nf}_2^2 +
\norm{f - S_Nf}_2^2 \geq \norm{S_Nf}_2^2 = \sum_{\abs n \leq N}
\abs{c_n}^2$; अब $N \to \infty$ लीजिए।
\end{proof}

\begin{example}[सर्वोत्तम सन्निकटन, मापा हुआ]\label{ex:b2:fourier:bestapprox}
द्विघात माध्य के अर्थ में निम्न घात के त्रिकोणमितीय बहुपद आरे-दाँत $f(t) = t$
($\intoo{-\pi}{\pi}$ पर) का कितना अच्छा सन्निकटन करते हैं? \cref{prop:b2:fourier:bessel} के अनुसार घात $N$ का
सर्वोत्तम सन्निकटन $S_N(f)$ \emph{ही} है, और वर्ग-त्रुटि
\[
\norm{f - S_Nf}_2^2 = \norm f_2^2 -
\sum_{\abs n\leq N}\abs{c_n}^2 .
\]
है। यहाँ $\norm f_2^2 = \frac{1}{2\pi}\int_{-\pi}^\pi t^2\dd t =
\frac{\pi^2}{3}$, और $b_n = \frac{2(-1)^{n+1}}{n}$ (\cref{ex:b2:fourier:basel}) से: $\abs{c_n}^2 + \abs{c_{-n}}^2 =
\frac{b_n^2}{2} = \frac{2}{n^2}$। अतः
\[
\norm{f - S_Nf}_2^2
= \frac{\pi^2}{3} - \sum_{n=1}^{N}\frac{2}{n^2}
\qquad\text{: संख्यात्मक रूप से } 1.29,\ 0.79,\ 0.57,\ 0.44
\]
$N = 1, 2, 3, 4$ के लिए --- अर्थात् घटती हुई, पर धीरे-धीरे: पुच्छ $\sum_{n>N}\frac2{n^2} \sim \frac2N$ गुणांकों के
धीमे $\frac1n$ क्षय से शासित है, और वह क्षय स्वयं उस उछाल का हस्ताक्षर है (\cref{exo:b2:fourier:6}
को उलटकर पढ़िए)। समापन दृष्टि: पारसेवाल सन्निकटन की गुणवत्ता को किसी
संख्यात्मक श्रेणी की पुच्छ में बदल देती है --- और किसी भी चित्र से पहले
बता देती है कि उछाल फूरिये श्रेणियों को अनिच्छा से अभिसरित कराते हैं।
\end{example}

\begin{method}[फूरिये गुणांक कुशलता से परिकलित करना]\label{met:b2:fourier:coefficients}
कुछ भी समाकलित करने से पहले:
\begin{enumerate}
  \item \emph{सम-विषमता:} सम $f$ के लिए $b_n = 0$, विषम $f$ के लिए $a_n
        = 0$, और बचे
        हुए समाकल $\frac2\pi
        \int_0^\pi$ तक सिमट जाते हैं --- आधा श्रम, दुगुनी विश्वसनीयता।
  \item \emph{त्रिकोणमितीय बहुपद पहले से ही हो चुके हैं:} गुणनफलों को
        रैखिक कीजिए ($\cos^3$, $\sin^2\cos$, \dots) और गुणांक सीधे पढ़ लीजिए (\cref{exo:b2:fourier:9});
        प्रसामान्य लांबिकता किसी और समाकलन को अनावश्यक बना देती है।
  \item \emph{चरघातांकियों के लिए सम्मिश्र चरघातांकी:} गुणनखंड $\eu^{at}$ अथवा
        अवमंदित दोलनों के लिए $c_n$ सीधे परिकलित कीजिए --- $\eu^{(a - \iu
        n)t}$ का एक ही
        समाकल दो खंडशः समाकलनों को हरा देता है (\cref{exo:b2:fourier:10})।
  \item \emph{किसी ज्ञात प्रसार का अवकलन:} यदि $f'$ के गुणांक ज्ञात हों
        और $f$ \omterm{def:b2:metric:continuity}{संतत} हो, तो $c_n(f) =
        \frac{c_n(f')}{\iu n}$ ($n \neq 0$) $c_0$ को छोड़कर सब कुछ पुनः दे
        देता है, और वह तो माध्य ही है --- यह प्रायः सबसे तेज़ रास्ता
        है, और ठीक \cref{thm:b2:fourier:parseval}~(1) की परिकल्पनाओं के अंतर्गत वैध।
\end{enumerate}
\end{method}

\section{डिरिक्ले की प्रमेय}

\begin{lemma}[डिरिक्ले कर्नेल]\label{lem:b2:fourier:kernel}
$S_N(f)(x) = \frac{1}{2\pi}\int_{-\pi}^{\pi} f(x + u)\,D_N(u)\,\dd
u$, जहाँ\index{डिरिक्ले कर्नेल}
\[
D_N(u) = \sum_{n=-N}^{N} \eu^{\iu nu}
= \frac{\sin\bigl((N + \frac12)u\bigr)}{\sin\frac u2}
\quad (u \notin 2\pi\Z),
\qquad
\frac{1}{2\pi}\int_{-\pi}^{\pi} D_N = 1 .
\]
\end{lemma}

\begin{proof}
$S_N$ में $c_n$ की परिभाषा डालिए और योग तथा समाकल आपस में बदल दीजिए (वैध: योग
परिमित है):
\[
S_N(f)(x)
= \sum_{n=-N}^{N}\Bigl(\frac{1}{2\pi}\int_{-\pi}^{\pi}
f(t)\,\eu^{-\iu nt}\dd t\Bigr)\eu^{\iu nx}
= \frac{1}{2\pi}\int_{-\pi}^{\pi} f(t)
\sum_{n=-N}^{N}\eu^{\iu n(x-t)}\,\dd t ;
\]
$u = t - x$ प्रतिस्थापित कीजिए और समाकल्य की $2\pi$-आवर्तिता से समाकलन-खंड को वापस
$\intcc{-\pi}{\pi}$ पर सरका दीजिए; और सममित सूचकांक-परिसर $\sum_n\eu^{-\iu nu} = D_N(u)$ बना देता है। बंद रूप:
अनुपात $\eu^{\iu u}$ वाला गुणोत्तर योग,
\[
D_N(u) = \eu^{-\iu Nu}\,\frac{\eu^{\iu(2N+1)u} - 1}{\eu^{\iu u} -
1}
= \frac{\eu^{\iu(N + \frac12)u} - \eu^{-\iu(N+\frac12)u}}
{\eu^{\iu u/2} - \eu^{-\iu u/2}} ,
\]
जो वही ज्या-भागफल है। और उसका माध्य $1$ है: केवल $n = 0$ योगदान देता है।
\end{proof}

\begin{theorem}[रीमान--लेबेग प्रमेयिका]\label{thm:b2:fourier:riemannlebesgue}
किसी खंड पर खंडशः \omterm{def:b2:metric:continuity}{संतत} $g$ के लिए $\lambda
\to +\infty$ होने पर $\int_a^b g(t)\sin(\lambda t + \varphi)\,\dd t \to 0$।\index{रीमान--लेबेग
प्रमेयिका}
\end{theorem}

\begin{proof}
$g$ का सोपान फलनों से एकसमान सन्निकटन कीजिए (\cref{thm:b2:funcseq:weierstrass} की आवश्यकता नहीं ---
खंडशः \omterm{def:b2:metric:continuity}{संतत} फलनों का प्रारंभिक सोपान-फलन सन्निकटन पर्याप्त है) और हर सोपान
का स्पष्ट समाकलन कीजिए: हर टुकड़ा $O\bigl(\frac1\lambda\bigr)$ का योगदान देता है, और सन्निकटन
त्रुटि $\varepsilon(b - a)$ का। यह तर्क प्रथम वर्ष के खंड के समाकलन अध्याय के अंतिम अभ्यास
के रूप में पूरा किया जा चुका है; और $C^1$ टुकड़ों के लिए इसके बदले खंडशः
समाकलन करके $\frac C\lambda$ से परिबद्ध किया जा सकता है।
\end{proof}

\begin{example}[गुणांक कितनी तेज़ी से मरते हैं?]\label{ex:b2:fourier:decaytable}
रीमान--लेबेग कहती है कि गुणांक $0$ की ओर जाते हैं; और उनकी \emph{दर}
चिकनाई का मापक है। इस अध्याय तथा उसके अभ्यासों से तीन नमूने:
\[
\text{वर्ग तरंग: } b_n = \frac{4}{\pi n}\ (n\ \text{विषम}),
\qquad
\abs t : a_n = \frac{-4}{\pi n^2}\ (n\ \text{विषम}),
\qquad
\abs{\sin t} : a_{2k} = \frac{-4}{\pi(4k^2-1)} .
\]
$f$ में कोई उछाल (वर्ग तरंग, आरे-दाँत) कोटि $\frac1n$ के गुणांक छोड़ जाता है:
अतः न \omterm{def:b2:funcseq:series}{सामान्य अभिसरण}, और उछालों पर गिब्स अतिक्रमण। और संततता के साथ कोई
कोना --- अर्थात् केवल $f'$ में उछाल --- कोटि को $\frac{1}{n^2}$ तक सुधार देता है: अतः
\omterm{def:b2:funcseq:series}{सामान्य अभिसरण}, एकसमान पुनर्निर्माण। सामान्य रूप में $k$ अवकलज $c_n =
O(n^{-k})$ (\cref{exo:b2:fourier:6})
ख़रीद लेते हैं, और विलोमतः हर घात से तेज़ क्षय वाला \omterm{def:b2:reduction:eigen}{वर्णक्रम} $f$ को $C^\infty$
होने पर बाध्य कर देता है (अब वैध रूप से पद-दर-पद अवकलन कीजिए)। समापन
दृष्टि: संकेत की नियमितता और \omterm{def:b2:reduction:eigen}{वर्णक्रम} का क्षय एक ही सूचना हैं --- कोई
अभियंता बिना फलन का आलेख बनाए एक को दूसरे की ढाल से पढ़ लेता है।
\end{example}

\begin{theorem}[डिरिक्ले]\label{thm:b2:fourier:dirichlet}
मान लीजिए $f$ $2\pi$-आवर्ती और खंडशः $C^1$ है। तब प्रत्येक $x$ के लिए
\[
S_N(f)(x) \xrightarrow[N \to \infty]{}
\frac{f(x^+) + f(x^-)}{2}
\]
(अर्थात् एकपक्षीय सीमाओं का माध्य) --- और विशेष रूप से हर संततता-बिंदु पर
$S_N(f)(x) \to
f(x)$।\index{डिरिक्ले की प्रमेय}
\end{theorem}


\begin{proof}
कर्नेल प्रमेयिका और उसके इकाई माध्य से, समाकल को $u > 0$ तथा $u < 0$ आधों में
बाँटने पर (हर एक का माध्य $\frac12$):
\begin{align*}
S_N(f)(x) - \frac{f(x^+) + f(x^-)}{2}
&= \frac{1}{2\pi}\int_{0}^{\pi} \bigl(f(x+u) -
f(x^+)\bigr)D_N(u)\,\dd u\\
&\quad+ \frac{1}{2\pi}\int_{-\pi}^{0}\bigl(f(x+u) -
f(x^-)\bigr)D_N(u)\,\dd u .
\end{align*}
पहले को निपटाइए (दूसरा \omterm{def:b2:quadratic:adjoint}{सममित} है)। लिखिए
\[
\bigl(f(x + u) - f(x^+)\bigr)\,D_N(u)
= \underbrace{\frac{f(x+u) - f(x^+)}{\sin\frac u2}}_{g(u)}\,
\sin\Bigl(\Bigl(N + \frac12\Bigr)u\Bigr) .
\]
फलन $g$ $\intoc{0}{\pi}$ पर खंडशः \omterm{def:b2:metric:continuity}{संतत} है और उसकी \emph{$0^+$ पर परिमित सीमा है}:
क्योंकि
\[
g(u) = \frac{f(x+u) - f(x^+)}{u}\cdot\frac{u}{\sin\frac u2} ,
\]
लिखने पर पहला गुणनखंड $f'(x^+)$ की ओर जाता है (खंडशः $C^1$ होने से एकपक्षीय
अवकलनीयता) और दूसरा $2$ की ओर ($v = \frac
u2$ पर मानक सीमा $\frac{\sin v}{v} \to 1$): अतः $g(0^+) = 2f'(x^+)$
विद्यमान है। इसलिए $g$ $\intcc{0}{\pi}$ तक खंडशः संततता के साथ बढ़ जाता है, और
रीमान--लेबेग (\cref{thm:b2:fourier:riemannlebesgue}) समाकल को $0$ पर भेज देती है। परिकल्पना का पूरा
उद्देश्य यही है: एकपक्षीय अवकलजों के बिना गुणनखंड $\frac{1}{\sin(u/2)}$ $0$ पर इतनी तेज़ी
से फट जाता है कि रीमान--लेबेग उसकी भरपाई नहीं कर सकती, और केवल \omterm{def:b2:metric:continuity}{संतत} $f$
के लिए \omterm{def:b2:funcseq:def}{बिंदुवार अभिसरण} सचमुच विफल हो सकता है --- और यही वह दरार है जिसे
फेयेर की प्रमेय (सप्ताहांत समस्या) औसत लेकर पाट देती है।
\end{proof}

\begin{example}[किसी उछाल पर डिरिक्ले]\label{ex:b2:fourier:jumpvalue}
$\intoo{-\pi}{\pi}$ पर आरे-दाँत $f(t) = t$ (\cref{ex:b2:fourier:basel} नीचे) के लिए आवर्ती विस्तार $t = \pi$ पर $f(\pi^-) = \pi$ से $f(\pi^+) = -\pi$
तक उछलता है। डिरिक्ले वहाँ मान $\frac{\pi + (-\pi)}{2} = 0$ का वचन देती है, और वास्तव में $\sum
\frac{2(-1)^{n+1}}{n}\sin nt$ का
हर पद $t = \pi$ पर लुप्त होता है: श्रेणी शालीनता से मध्यबिंदु पर अभिसरित होती
है और दोनों एकपक्षीय मानों की उपेक्षा कर देती है। और मूल्यांकन-बिंदु को
$t =
\frac\pi2$ (कोई संततता-बिंदु) पर ले जाने से वही श्रेणी लाइब्नित्स की $\frac\pi4$ बन
जाती है। एक ही श्रेणी, दो व्यवहार --- अर्थात् ठीक प्रमेय के दोनों उपवाक्य।
\end{example}

\begin{theorem}[$C^1$ के लिए सामान्य अभिसरण; पारसेवाल]\label{thm:b2:fourier:parseval}
\begin{enumerate}
  \item यदि $f$ \omterm{def:b2:metric:continuity}{संतत}, $2\pi$-आवर्ती और खंडशः $C^1$ हो, तो $c_n(f') = \iu n\,c_n(f)$, $f$ की फूरिये
        श्रेणी $\R$ पर \emph{सामान्य} रूप से अभिसरित होती है, और उसका योग
        $f$ है।
  \item (पारसेवाल) प्रत्येक खंडशः \omterm{def:b2:metric:continuity}{संतत} $2\pi$-आवर्ती $f$ के
        लिए:\index{पारसेवाल की सर्वसमिका}
        \[
        \frac{1}{2\pi}\int_{-\pi}^{\pi}\abs{f}^2
        = \sum_{n=-\infty}^{\infty} \abs{c_n(f)}^2
        = \frac{\abs{a_0}^2}{4} + \frac12\sum_{n\geq1}
        \bigl(\abs{a_n}^2 + \abs{b_n}^2\bigr).
        \]
        \emph{(यहाँ $f$ के \omterm{def:b2:metric:continuity}{संतत} तथा खंडशः $C^1$ होने पर सिद्ध; सामान्य
        स्थिति में स्वीकृत।)}
\end{enumerate}
\end{theorem}

\begin{proof}
(1) हर $C^1$ टुकड़े पर खंडशः समाकलन (सीमा-पद संततता और आवर्तिता से कट जाते
हैं): $c_n(f') = \iu n c_n(f)$। फिर दोनों वर्ग-योग्य कुलों पर कोशी--श्वार्ज़ ($f'$ के लिए \cref{prop:b2:fourier:bessel}):
\[
\sum_{n \neq 0} \abs{c_n(f)} = \sum_{n\neq0}
\frac{\abs{c_n(f')}}{\abs n}
\leq \Bigl(\sum \abs{c_n(f')}^2\Bigr)^{1/2}
\Bigl(\sum_{n\neq0}\frac{1}{n^2}\Bigr)^{1/2} < \infty :
\]
अर्थात् फूरिये श्रेणी का \omterm{def:b2:funcseq:series}{सामान्य अभिसरण}। उसका योग \omterm{def:b2:metric:continuity}{संतत} है और डिरिक्ले
(\cref{thm:b2:fourier:dirichlet}: $f$ \omterm{def:b2:metric:continuity}{संतत} है) से हर बिंदु पर $f$ से मेल खाता है: अतः श्रेणी $f$ पर
\omterm{def:b2:funcseq:def}{एकसमान रूप से} अभिसरित होती है।

(2) ऐसे $f$ के लिए: \omterm{def:b2:funcseq:def}{एकसमान रूप से} $S_N f \to f$, अतः $\norm{f - S_Nf}_2
\leq \norm{f - S_Nf}_\infty \to 0$, और पाइथागोरस ($\norm f_2^2 = \sum_{\abs n \leq N}\abs{c_n}^2 + \norm{f -
S_Nf}_2^2$) सीमा
तक चला जाता है। वास्तविक रूप $a_n, b_n$ के साथ केवल लेखा-जोखा है।
\end{proof}

\begin{example}[$C^1$ पुच्छ-परिबंध, परिमाणात्मक रूप में]\label{ex:b2:fourier:tailbound}
\cref{thm:b2:fourier:parseval}~(1) की उपपत्ति में एक काम आने वाला आकलन छिपा है। \omterm{def:b2:metric:continuity}{संतत} खंडशः-$C^1$ $f$
के लिए वही कोशी--श्वार्ज़ केवल पुच्छ पर लगाने से
\[
\sum_{\abs n > N}\abs{c_n(f)}
= \sum_{\abs n > N}\frac{\abs{c_n(f')}}{\abs n}
\leq \Bigl(\sum_{\abs n > N}\abs{c_n(f')}^2\Bigr)^{\!1/2}
\Bigl(\sum_{\abs n>N}\frac{1}{n^2}\Bigr)^{\!1/2}
\leq \norm{f'}_2\,\sqrt{\frac{2}{N}} ,
\]
मिलता है, जिसमें $f'$ के लिए बेसेल और $\sum_{n>N}n^{-2} \leq \frac1N$ का उपयोग हुआ। अतः आंशिक योगों की
एकसमान त्रुटि
\[
\norm{f - S_Nf}_\infty
\leq \sum_{\abs n>N}\abs{c_n(f)}
\leq \norm{f'}_2\,\sqrt{\frac2N} .
\]
का पालन करती है। $f(t) = \abs t$ के लिए: $\norm{f'}_2 = 1$ (अवकलज $\pm1$ है), अतः दस पद ही $\abs t$ को
\omterm{def:b2:funcseq:def}{एकसमान रूप से} $\sqrt{0.2} \approx 0.45$ के भीतर पुनर्निर्मित कर देते हैं, और $N = 10^4$ $0.015$ के
भीतर। समापन दृष्टि: एक अवकलज $\frac{1}{\sqrt N}$ वाली एकसमान दर ख़रीद लेता है; और वर्ग
तरंग की $\frac1n$-गुणांक वाली दुनिया (जहाँ \omterm{def:b2:funcseq:def}{एकसमान अभिसरण} है ही नहीं) से तुलना
करने पर \cref{ex:b2:fourier:decaytable} के शब्दकोश को संख्याएँ मिल जाती हैं।
\end{example}

\begin{example}[बाज़ल और साथी]\label{ex:b2:fourier:basel}
$\intoo{-\pi}{\pi}$ पर $f(t) = t$ लीजिए, और उसे $2\pi$-आवर्ती रूप से बढ़ा दीजिए (आरे-दाँत, खंडशः
$C^1$)। परिकलन करने पर $a_n = 0$ (विषमता से) और
\[
b_n = \frac{1}{\pi}\int_{-\pi}^{\pi} t\sin nt\,\dd t
= \frac{2(-1)^{n+1}}{n} .
\]
$t = \frac\pi2$ पर डिरिक्ले लाइब्नित्स का $\frac\pi4 = 1 -
\frac13 + \frac15 - \dots$ लौटा देती है; और पारसेवाल देती है
\[
\frac{1}{2\pi}\int_{-\pi}^{\pi} t^2\,\dd t = \frac{\pi^2}{3}
= \frac12\sum_{n\geq1}\frac{4}{n^2}
\quad\Longrightarrow\quad
\boxed{\;\sum_{n\geq1}\frac{1}{n^2} = \frac{\pi^2}{6}\;}
\]
--- अर्थात् ऑयलर का बाज़ल योग, दो पंक्तियों में। फलन $f(t) = t^2$ इसी तरह $\sum \frac1{n^4} = \frac{\pi^4}{90}$
देता है (\cref{exo:b2:fourier:3})।
\end{example}

\begin{example}[अंतर्निहित जाँच के साथ एक पूरा प्रसार: $\abs{\sin t}$]\label{ex:b2:fourier:abssin}
फलन $f(t) = \abs{\sin t}$ \omterm{def:b2:metric:continuity}{संतत}, सम, $\pi$-आवर्ती (अतः $2\pi$-आवर्ती) और खंडशः $C^1$ है। समता $b_n$
मार देती है; $a_0 = \frac1\pi\int_{-\pi}^{\pi}
\abs{\sin t}\dd t = \frac4\pi$; और $n \geq 1$ के लिए गुणनफल-से-योग देता है
\[
a_n = \frac2\pi\int_0^\pi \sin t\cos nt\,\dd t
= \frac{1}{\pi}\int_0^\pi\bigl(\sin(1+n)t +
\sin(1-n)t\bigr)\dd t
= \frac2\pi\cdot\frac{1 + \cos n\pi}{1 - n^2}
\]
$n \neq 1$ के लिए (और $a_1 = 0$ सीधे): विषम $n$ के लिए शून्य, और $a_{2k} = \frac{-4}{\pi(4k^2-1)}$। \cref{thm:b2:fourier:parseval}~(1) के
अनुसार अभिसरण सामान्य है, और
\[
\abs{\sin t} = \frac{2}{\pi} - \frac{4}{\pi}
\sum_{k\geq1}\frac{\cos(2kt)}{4k^2 - 1}
\qquad (t \in \R) .
\]
$t = 0$ पर अंतर्निहित जाँच: सर्वसमिका $\sum_{k\geq1}\frac{1}{4k^2-1} = \frac12$ माँगती है, जिसकी पुष्टि क्रमिक
निरसन कर देता है:
\[
\sum_{k\geq1}\frac{1}{4k^2-1}
= \frac12\sum_{k\geq1}\Bigl(\frac{1}{2k-1} -
\frac{1}{2k+1}\Bigr) = \frac12 . \checkmark
\]
समापन दृष्टि: $\abs{\sin}$ का \omterm{def:b2:reduction:eigen}{वर्णक्रम} केवल \emph{सम} आवृत्तियों पर बसता है ---
किसी ज्या को दिष्ट करने से उसकी आवृत्ति-सामग्री दुगुनी हो जाती है, और
इसीलिए पूर्ण-तरंग दिष्टकारी $100$ या $120$ हर्ट्ज़ पर गूँजते हैं, यानी मुख्य
आवृत्ति से दुगुने पर।
\end{example}

\begin{example}[परिकलन के यंत्र के रूप में पारसेवाल]\label{ex:b2:fourier:zetafourodd}
पारसेवाल प्रसारों को थोक में संख्यात्मक श्रेणियों में बदल देती है। उसे
$f(t) = \abs t$ पर लगाइए (\cref{exo:b2:fourier:2}: $a_0 =
\pi$, विषम $n$ के लिए $a_n = \frac{-4}{\pi n^2}$, शेष शून्य):
\[
\frac{1}{2\pi}\int_{-\pi}^{\pi}t^2\,\dd t = \frac{\pi^2}{3}
= \frac{a_0^2}{4} + \frac12\sum_{n \text{ विषम}} a_n^2
= \frac{\pi^2}{4} +
\frac{8}{\pi^2}\sum_{n\text{ विषम}}\frac{1}{n^4} ,
\]
जिससे
\[
\sum_{n\text{ विषम}}\frac{1}{n^4}
= \frac{\pi^2}{8}\Bigl(\frac{\pi^2}{3} -
\frac{\pi^2}{4}\Bigr) = \frac{\pi^4}{96} .
\]
\cref{exo:b2:fourier:3} के विरुद्ध प्रतिजाँच: $\sum\frac1{n^4}$ को विषम और सम भागों में बाँटने पर
$\zeta$-प्रकार का लेखा-जोखा $S = S_{\mathrm{odd}} + \frac{S}{16}$ मिलता है, अतः $S =
\frac{16}{15}\cdot\frac{\pi^4}{96} = \frac{\pi^4}{90}$ --- ठीक वही मान जो वहाँ
किसी दूसरे फलन से मिला था। दो प्रसार, एक संख्या: और यह संगति पारसेवाल की
समदूरिकता का काम है। समापन दृष्टि: हर नया फूरिये प्रसार श्रेणी-सर्वसमिकाएँ
उत्पन्न करने वाली मशीन है; और सप्ताहांत समस्या का भाग II समझाता है कि यह
मशीन स्वयं का खंडन कभी क्यों नहीं कर सकती।
\end{example}

\begin{omfigure}
\begin{tikzpicture}
\begin{axis}[omaxis, width=9.2cm, height=5.6cm,
    xmin=-3.6, xmax=3.6, ymin=-1.5, ymax=1.5,
    xtick={-3.1416, 0, 3.1416},
    xticklabels={$-\pi$, $0$, $\pi$}, ytick={-1,1}]
  \addplot[omExo, thick] coordinates {(-3.1416,-1) (0,-1)};
  \addplot[omExo, thick] coordinates {(0,1) (3.1416,1)};
  \addplot[omDef, thick, domain=-3.3:3.3, samples=200]
    {4/pi*sin(deg(x))};
  \addplot[omThm, very thick, domain=-3.3:3.3, samples=400]
    {4/pi*(sin(deg(x)) + sin(deg(3*x))/3 + sin(deg(5*x))/5 +
     sin(deg(7*x))/7 + sin(deg(9*x))/9)};
\end{axis}
\end{tikzpicture}

{\small वर्ग तरंग (धूसर) और \omterm{def:b2:fourier:coefficients}{फूरिये आंशिक योग} $S_1$ (नीला) तथा $S_9$ (लाल):
हर संततता-बिंदु पर अभिसरण, परंतु उछालों के पास टिका हुआ $\sim 9\%$ अतिक्रमण ---
अर्थात् \emph{गिब्स परिघटना}\index{गिब्स परिघटना}। \omterm{def:b2:funcseq:def}{एकसमान अभिसरण} ठीक
इसलिए विफल है कि सीमा असंतत है।}
\end{omfigure}

\begin{example}[स्थानांतरण और मॉडुलन]\label{ex:b2:fourier:shiftrules}
एक-एक पंक्ति के दो नियम एक ही प्रसार से अनेक प्रसार उत्पन्न कर देते हैं।
$a
\in \R$ के लिए $s = t - a$ प्रतिस्थापित करने पर:
\[
c_n\bigl(f(\cdot - a)\bigr)
= \frac{1}{2\pi}\int_{-\pi}^{\pi}f(t - a)\eu^{-\iu nt}\dd t
= \eu^{-\iu na}\,c_n(f)
\qquad\text{(स्थानांतरण वर्णक्रम का मॉडुलन करता है)},
\]
और सीधे परिभाषा से,
\[
c_n\bigl(\eu^{\iu kt}f\bigr) = c_{n-k}(f)
\qquad\text{(मॉडुलन वर्णक्रम का स्थानांतरण करता है)}.
\]
किया हुआ उदाहरण: $\pi$ से खिसकाया गया आरे-दाँत, अर्थात् $f(t) = t$ के साथ $g(t) = f(t -
\pi)$,
जिसके $b_n$-गुणांक $(-1)^n\cdot\frac{2(-1)^{n+1}}{n} = -\frac2n$ हैं: अर्थात् $0$ के बदले $\pi$ पर उछलने वाले
आरे-दाँत का प्रसार $g \sim -2\sum\frac{\sin nt}{n}$ --- और कोई समाकल फिर से नहीं करना पड़ा। समापन
दृष्टि: समय के खिसकाव केवल कला बदलते हैं, आयाम कभी नहीं ($\abs{c_n}$
खिसकाव-अपरिवर्त्य है), और इसीलिए ऊर्जा (पारसेवाल) तथा अभिसरण-वर्ग संकेत के
\emph{आकार} के गुणधर्म हैं, इस बात के नहीं कि घड़ी कहाँ से चलनी शुरू हुई।
\end{example}

\begin{remark}[सामान्य भूलें]
\emph{(क) तीन अभिसरण, तीन मुद्राएँ:} \omterm{def:b2:funcseq:def}{बिंदुवार} (डिरिक्ले: खंडशः $C^1$
चाहिए, और हर उछाल पर \emph{मध्यबिंदु} देती है --- कभी एकपक्षीय मान नहीं),
एकसमान (\omterm{def:b2:metric:continuity}{संतत} सीमा चाहिए; उछाल के पार असंभव, और गिब्स उसका दिखाई देने
वाला लक्षण है), और द्विघात माध्य (पारसेवाल: सबसे मज़बूत, और अलग-अलग
बिंदुओं के प्रति अंधा)। सदा बताइए कि आप कौन-सा दावा कर रहे हैं। \emph{(ख)
अपने आप पद-दर-पद अवकलन नहीं:} \cref{ex:b2:fourier:basel} की आरे-दाँत श्रेणी का पद-दर-पद अवकलन
$\sum
2(-1)^{n+1}\cos nt$ देता है, जिसके पद $0$ की ओर तक नहीं जाते --- फलन-अनुक्रम वाले अध्याय
की स्थानांतरण प्रमेयों को \emph{व्युत्पन्न} श्रेणी का \omterm{def:b2:funcseq:def}{एकसमान अभिसरण} चाहिए,
जिसे उछाल नष्ट कर देता है। पहले चिकनाई, फिर अवकलन (\cref{exo:b2:fourier:6} शब्दकोश है)।
\emph{(ग) सममित आंशिक योग:} डिरिक्ले की प्रमेय $S_N =
\sum_{-N}^{N}$ के बारे में है; और
पुनर्विन्यास अथवा पहले एक ही ओर का योग लेने से अपसरण अभिसरण में और वापस
बदल सकता है। \emph{(घ) सामान्यीकरण का बहाव:} परिपाटियाँ पुस्तकों में
भिन्न होती हैं (आगे $\frac{1}{2\pi}$ अथवा $\frac1\pi$, आवर्तकाल $2\pi$ अथवा $1$); भरोसेमंद
अपरिवर्त्य प्रसामान्य लांबिकता के संबंध हैं --- किसी भी सूत्र पर भरोसा
करने से पहले उपस्थित परिपाटी में $\langle e_m, e_n\rangle$ फिर से परिकलित कीजिए।
\end{remark}

\begin{remark}[यह कहाँ काम आता है]
पारसेवाल फूरिये श्रेणियों के $L^2$ सिद्धांत का बीज है: तृतीय वर्ष का खंड
चित्र पूरा कर देता है (चरघातांकी $L^2$ का हिल्बर्ट आधार हैं, और प्रतिचित्रण
$f \mapsto (c_n)$ एकैकी आच्छादक समदूरिक है)। इसी खंड के भीतर सप्ताहांत समस्या
\emph{फेयेर की प्रमेय} सिद्ध करती है --- हर \omterm{def:b2:metric:continuity}{संतत} आवर्ती $f$ के लिए फूरिये
श्रेणी के चेज़ारो माध्य \omterm{def:b2:funcseq:def}{एकसमान रूप से} अभिसरित होते हैं --- जो स्वीकृत
सामान्य पारसेवाल को प्रमेय बना देती है, त्रिकोणमितीय वाइरश्ट्रास प्रमेय
देती है, और दो शानदार लाभ चुकाती है: वाइल की समवितरण प्रमेय और समपरिमापी
असमिका। अनुप्रयुक्त गणित इस अध्याय को रोज़ पढ़ता है: संकेतों के \omterm{def:b2:reduction:eigen}{वर्णक्रम},
कंपित निकायों की संनादियाँ, और तीव्र फूरिये रूपांतरण (जिसका परिमित अवतार
\cref{exo:b2:hermitian:10} था)।
\end{remark}

\begin{remark}[इसी खंड के भीतर के परिप्रेक्ष्य]
तीन अध्याय इससे बातचीत करते हैं। पीछे की ओर: \omterm{def:b2:hermitian:adjoint}{हर्मीशियन} अध्याय ने ज्यामिति
दी (प्रसामान्य लांबिक कुल, प्रक्षेप, बेसेल), और फलन-अनुक्रम अध्याय ने
विश्लेषण (\omterm{def:b2:funcseq:def}{एकसमान अभिसरण}, स्थानांतरण प्रमेयें, सन्निकट तत्समक --- फेयेर
कर्नेल फूरिये श्रेणियों के लिए वही है जो \omterm{thm:b2:funcseq:weierstrass}{बर्नस्टाइन बहुपद} वाइरश्ट्रास के
लिए थे)। बग़ल में: घात-श्रेणी अध्याय का सीमा-सिद्धांत \omterm{thm:b2:series:abel}{आबेल योग} के द्वारा
लौटता है, जिसे यहाँ प्वासों कर्नेल साकार करता है (\cref{exo:b2:fourier:12}) --- जहाँ चक्रिका की
त्रिज्या $r$ योग-प्राचल की भूमिका निभाती है। और आगे की ओर: अवकल समीकरणों
वाला अध्याय आवर्ती प्रणोदन को संनादियों में तोड़कर हर एक को दोलक की
आवृत्ति-प्रतिक्रिया में डालता है; अनुनाद तब होता है जब निवेश की कोई फूरिये
विधा किसी प्राकृतिक आवृत्ति से मेल खा जाए, और इसीलिए उसकी सप्ताहांत समस्या
और इस अध्याय की समस्या एक ही कहानी के दो आधे हैं।
\end{remark}

\begin{omfigure}
\begin{tikzpicture}
\begin{axis}[omaxis, width=9.2cm, height=5.6cm,
    xmin=-3.6, xmax=3.6, ymin=-2.5, ymax=9.5,
    xtick={-3.1416, 0, 3.1416},
    xticklabels={$-\pi$, $0$, $\pi$}, ytick={8}]
  \addplot[omDef, thick, domain=0.02:3.3, samples=300]
    {sin(deg(8.5*x))/sin(deg(0.5*x))};
  \addplot[omDef, thick, domain=-3.3:-0.02, samples=300]
    {sin(deg(8.5*x))/sin(deg(0.5*x))};
  \addplot[omThm, very thick, domain=0.02:3.3, samples=300]
    {sin(deg(4*x))^2/(8*sin(deg(0.5*x))^2)};
  \addplot[omThm, very thick, domain=-3.3:-0.02, samples=300]
    {sin(deg(4*x))^2/(8*sin(deg(0.5*x))^2)};
  \node[font=\small, omDef] at (axis cs:1.05,5.2) {$D_8$};
  \node[font=\small, omThm] at (axis cs:-0.75,6.2) {$F_8$};
\end{axis}
\end{tikzpicture}

{\small \omterm{lem:b2:fourier:kernel}{डिरिक्ले कर्नेल} $D_8$ (नीला) दोलन करता है और ऋणात्मक मान लेता है;
जबकि फेयेर कर्नेल $F_8$ (लाल) अऋणात्मक है, $0$ पर संकेंद्रित होता है, और
उसका माध्य $1$ है: अर्थात् एक \emph{सन्निकट तत्समक}। \omterm{lem:b2:fourier:kernel}{डिरिक्ले कर्नेल} में
ठीक धनात्मकता की कमी है, और वही फेयेर की प्रमेय को सप्ताहांत समस्या में
अशर्त बना देती है।}
\end{omfigure}

\section{अभ्यास}

\begin{exercise}[$\star$]\label{exo:b2:fourier:1}
वर्ग तरंग ($\intoo{-\pi}{0}$ पर $f = -1$, $\intoo{0}{\pi}$ पर $+1$) के \omterm{def:b2:fourier:coefficients}{फूरिये गुणांक} परिकलित कीजिए, $t = \frac\pi2$
पर तथा उछाल $t = 0$ पर डिरिक्ले का निष्कर्ष बताइए, और लाइब्नित्स की श्रेणी
पुनः प्राप्त कीजिए।
\end{exercise}

\begin{exercise}[$\star$]\label{exo:b2:fourier:2}
$f(t) = \abs t$ ($\abs t \leq \pi$, $2\pi$-आवर्ती) का फूरिये श्रेणी में प्रसार कीजिए; \emph{सामान्य}
अभिसरण न्यायोचित ठहराइए; $t =
0$ पर मूल्यांकन करके $\sum_{k\geq0}\frac{1}{(2k+1)^2} = \frac{\pi^2}{8}$ प्राप्त कीजिए, और उससे
बाज़ल पुनः निकालिए।
\end{exercise}

\begin{exercise}[$\star$]\label{exo:b2:fourier:3}
$f(t) = t^2$ ($\abs t \leq \pi$) का प्रसार कीजिए और निष्कर्ष रूप में निकालिए
\[
\sum_{n\geq1}\frac{(-1)^{n+1}}{n^2} = \frac{\pi^2}{12},
\qquad
\sum_{n\geq1}\frac{1}{n^4} = \frac{\pi^4}{90}
\quad\text{(पारसेवाल)}.
\]
\end{exercise}

\begin{exercise}[$\star\star$]\label{exo:b2:fourier:4}
मान लीजिए $f$ \omterm{def:b2:metric:continuity}{संतत} $2\pi$-आवर्ती है और सभी $n$ के लिए $c_n(f) = 0$। सिद्ध कीजिए कि
$f = 0$ \emph{(पारसेवाल --- यहाँ वह किस वर्ग के लिए सिद्ध हुई है? न्यायोचित
ठहराइए कि सामान्य पारसेवाल स्वीकार कर लेने पर संततता तथा खंडशः $C^1$ की शर्त
हटाई जा सकती है, अथवा सघनता वाला तर्क रूपरेखा में दीजिए)}।
\end{exercise}

\begin{exercise}[$\star\star$]\label{exo:b2:fourier:5}
$\alpha \notin \Z$ के लिए $f(t) = \cos(\alpha t)$ ($\abs t
\leq \pi$) का प्रसार कीजिए और कोस्पर्शज्या का आंशिक-भिन्न प्रसार
निष्कर्ष रूप में निकालिए:
\[
\pi\cot(\pi\alpha) = \frac{1}{\alpha} + \sum_{n\geq1}
\frac{2\alpha}{\alpha^2 - n^2} .
\]
\end{exercise}

\begin{exercise}[$\star\star$]\label{exo:b2:fourier:6}
सिद्ध कीजिए कि यदि $f$ $2\pi$-आवर्ती और $C^k$ हो, जहाँ $f^{(k)}$ खंडशः \omterm{def:b2:metric:continuity}{संतत} है, तो
$c_n(f) = O\bigl(\abs n^{-k}\bigr)$: अर्थात् संकेत की चिकनाई $=$ उसके \omterm{def:b2:reduction:eigen}{वर्णक्रम} का क्षय।
\end{exercise}

\begin{exercise}[$\star\star\star$]\label{exo:b2:fourier:7}
(विर्टिंगर असमिका) मान लीजिए $f$ $C^1$, $2\pi$-आवर्ती है और $\int_{-\pi}^{\pi} f = 0$। सिद्ध कीजिए
\[
\int_{-\pi}^{\pi} \abs{f}^2 \leq \int_{-\pi}^{\pi} \abs{f'}^2 ,
\]
और समता तभी जब $f(t) = a\cos t + b\sin t$। \emph{(दोनों पक्षों पर पारसेवाल; $\abs{c_n}^2$ और $n^2\abs{c_n}^2$ की तुलना
कीजिए।)}
\end{exercise}

\begin{exercise}[$\star\star\star$]\label{exo:b2:fourier:8}
(गिब्स अचर) \cref{exo:b2:fourier:1} की वर्ग तरंग के लिए आंशिक योग का $x_N =
\frac{\pi}{2N}$ पर मूल्यांकन कीजिए:
$u_k = \frac{(2k+1)\pi}{2N}$ और $\Delta u
= \frac{\pi}{N}$ लिखकर दिखाइए कि
\[
S_{2N-1}\Bigl(\frac{\pi}{2N}\Bigr)
= \frac{4}{\pi}\sum_{k=0}^{N-1} \frac{\sin u_k}{2k+1}
= \frac{2}{\pi}\sum_{k=0}^{N-1} \frac{\sin u_k}{u_k}\,\Delta u
\xrightarrow[N\to\infty]{}
\frac{2}{\pi}\int_0^{\pi}\frac{\sin u}{u}\,\dd u \approx 1.179 :
\]
जो मध्यबिंदुओं पर $\intcc{0}{\pi}$ पर $\frac{2}{\pi}\cdot\frac{\sin u}{u}$ का रीमान योग है। निष्कर्ष निकालिए कि उछाल के
मान $1$ से आगे का अतिक्रमण $N \to \infty$ होने पर भी लुप्त नहीं होता।
\end{exercise}

\begin{exercise}[$\star$]\label{exo:b2:fourier:9}
$\cos^3 t$ और $\sin^2 t\,\cos t$ का फूरिये श्रेणी में प्रसार कीजिए \emph{(रैखिक कीजिए; गुणांकों
की अद्वितीयता से कोई त्रिकोणमितीय बहुपद अपनी ही फूरिये श्रेणी होता है)}।
हर एक के लिए $c_n$, $a_n$, $b_n$ क्या हैं?
\end{exercise}

\begin{exercise}[$\star\star$]\label{exo:b2:fourier:10}
मान लीजिए $a > 0$ और $\intoc{-\pi}{\pi}$ पर $f(t) = \eu^{at}$, जिसे $2\pi$-आवर्ती रूप से बढ़ाया गया है। \[
c_n(f) = \frac{(-1)^n\sinh(a\pi)}{\pi(a - \iu n)} ,
\]
परिकलित कीजिए, उछाल $t = \pi$ पर डिरिक्ले की प्रमेय लगाइए, और अतिपरवलयिक
कोस्पर्शज्या का आंशिक-भिन्न प्रसार निष्कर्ष रूप में निकालिए:
\[
\coth(\pi a) = \frac{1}{\pi a} + \sum_{n\geq1}
\frac{2a}{\pi(a^2 + n^2)} .
\]
\end{exercise}

\begin{exercise}[$\star\star$]\label{exo:b2:fourier:11}
(संवलन) \omterm{def:b2:metric:continuity}{संतत} $2\pi$-आवर्ती $f, g$ के लिए
\[
(f * g)(x) = \frac{1}{2\pi}\int_{-\pi}^{\pi} f(x - t)\,g(t)\,
\dd t .
\]
परिभाषित कीजिए। दिखाइए कि $f * g = g * f$, कि $c_n(f*g) = c_n(f)\,c_n(g)$ \emph{(किसी \omterm{def:b2:metric:continuity}{संतत} समाकल्य के दोनों
समाकल आपस में बदलने के लिए ऊपरी सीमा के दोनों फलनों की तुलना कीजिए: दोनों
बाएँ अंत्यबिंदु पर लुप्त हैं और संततता तथा समाकल-चिह्न के भीतर अवकलन से
उनके अवकलज एक ही हैं)}, और यह कि \omterm{lem:b2:fourier:kernel}{डिरिक्ले कर्नेल} के लिए $S_N(f) =
f * D_N$। (सप्ताहांत
समस्या के फेयेर माध्य भी इसी तरह संवलन हैं, $\sigma_N(f) = f *
F_N$।)
\end{exercise}

\begin{exercise}[$\star\star\star$]\label{exo:b2:fourier:12}
(प्वासों कर्नेल: फूरिये श्रेणियों के आबेल माध्य) $0 \leq r <
1$ के लिए $P_r(t) = \sum_{n\in\Z} r^{\abs n}\eu^{\iu nt}$ रखिए।
\begin{enumerate}
  \item दोनों गुणोत्तर श्रेणियों का योग लीजिए और दिखाइए
        \[
        P_r(t) = \frac{1 - r^2}{1 - 2r\cos t + r^2} > 0,
        \qquad
        \frac{1}{2\pi}\int_{-\pi}^{\pi}P_r = 1 .
        \]
  \item दिखाइए कि $\delta \leq \abs t \leq \pi$ के लिए $r \to 1^-$ होने पर \omterm{def:b2:funcseq:def}{एकसमान रूप से} $P_r(t)
        \leq \frac{1 - r^2}{1 - 2r\cos\delta + r^2} \to 0$।
  \item निष्कर्ष निकालिए कि प्रत्येक \omterm{def:b2:metric:continuity}{संतत} $2\pi$-आवर्ती $f$ के लिए $r \to 1^-$ होने
        पर \emph{आबेल माध्य} $(f * P_r)(x) = \sum_n
        r^{\abs n}c_n(f)\,\eu^{\iu nx}$ \omterm{def:b2:funcseq:def}{एकसमान रूप से} $f$ पर अभिसरित होते हैं
        --- अर्थात् फेयेर की प्रमेय का संतत-प्राचल वाला भाई, और घात-श्रेणी
        अध्याय के \omterm{thm:b2:series:abel}{आबेल योग} का फूरिये अवतार।
\end{enumerate}
\end{exercise}

\section{समस्या: फेयेर की प्रमेय और उसके लाभ}

\begin{problem}\label{pb:b2:fourier:1}
डिरिक्ले की प्रमेय को $f$ का खंडशः $C^1$ होना चाहिए; और केवल \omterm{def:b2:metric:continuity}{संतत} $f$ के
लिए आंशिक योग $S_N(f)$ बिगड़ सकते हैं। फेयेर की खोज: उनके \emph{चेज़ारो माध्य}
कभी नहीं बिगड़ते। इंजन फेयेर कर्नेल की धनात्मकता है, और फ़सल विशाल है:
एकसमान त्रिकोणमितीय सन्निकटन (वाइरश्ट्रास), फूरिये गुणांकों की
अद्वितीयता, इस अध्याय के हर फलन के लिए पारसेवाल (\cref{thm:b2:fourier:parseval} का ``स्वीकृत''
हटाते हुए), वाइल की समवितरण प्रमेय, और --- सदी भर की ज्यामिति पर मुकुट
रखते हुए --- समपरिमापी असमिका।\index{फेयेर की प्रमेय}\index{समवितरण}\index{समपरिमापी असमिका} आगे
सर्वत्र $f$ $2\pi$-आवर्ती और खंडशः \omterm{def:b2:metric:continuity}{संतत} है, तथा
\[
\sigma_N(f) = \frac{S_0(f) + S_1(f) + \dots +
S_{N-1}(f)}{N} .
\]

\medskip
\noindent\textbf{भाग I --- फेयेर कर्नेल।}
\begin{enumerate}
  \item दिखाइए कि $F_N = \frac{D_0 + \dots + D_{N-1}}{N}$ सहित $\sigma_N(f)(x) =
        \frac{1}{2\pi}\int_{-\pi}^{\pi}f(x+u)\,F_N(u)\,\dd u$, और यह कि $\frac{1}{2\pi}\int_{-\pi}^{\pi}F_N = 1$।
  \item $u \notin 2\pi\Z$ के लिए बंद रूप सिद्ध कीजिए:
        \[
        F_N(u) = \frac{1}{N}\,
        \frac{\sin^2\bigl(\frac{Nu}{2}\bigr)}
        {\sin^2\bigl(\frac u2\bigr)} \;\geq\; 0
        \]
        \emph{($\sin\bigl((n+\frac12)u\bigr)$ का योग किसी गुणोत्तर श्रेणी के काल्पनिक भाग के रूप में
        लीजिए)}।
  \item संकेंद्रण-आकलन दिखाइए: $0 < \delta \leq
        \abs u \leq \pi$ के लिए
        \[
        F_N(u) \leq \frac{1}{N\sin^2\frac\delta2}
        \xrightarrow[N\to\infty]{} 0
        \quad\text{एकसमान रूप से} :
        \]
        अर्थात् $(F_N)$ धनात्मक सन्निकट तत्समक है।
  \item (फेयेर की प्रमेय) सिद्ध कीजिए: यदि $f$ \omterm{def:b2:metric:continuity}{संतत} और $2\pi$-आवर्ती हो,
        तो $\R$ पर \emph{\omterm{def:b2:funcseq:def}{एकसमान रूप से}} $\sigma_N(f) \to f$ \emph{($\bigl(f(x+u) - f(x)\bigr)F_N(u)$ के समाकल को $\abs u =
        \delta$
        पर बाँटिए; हाइने तथा प्रश्न 1--3 का उपयोग कीजिए)}।
  \item खंडशः \omterm{def:b2:metric:continuity}{संतत} $f$ के लिए हर $x$ पर \omterm{def:b2:funcseq:def}{बिंदुवार} रूप $\sigma_N(f)(x) \to \frac{f(x^+) +
        f(x^-)}{2}$ दिखाइए, और
        एकसमान परिबंध
        $\norm{\sigma_N(f)}_\infty \leq \norm f_\infty$
        भी \emph{(धनात्मकता!)}।
\end{enumerate}

\medskip
\noindent\textbf{भाग II --- वाइरश्ट्रास, अद्वितीयता, पारसेवाल।}
\begin{enumerate}[resume]
  \item (त्रिकोणमितीय वाइरश्ट्रास) निष्कर्ष निकालिए: हर \omterm{def:b2:metric:continuity}{संतत} $2\pi$-आवर्ती
        फलन त्रिकोणमितीय बहुपदों की एकसमान सीमा है।
  \item (अद्वितीयता) निष्कर्ष निकालिए: जिस \omterm{def:b2:metric:continuity}{संतत} $f$ के लिए सभी $n$ हेतु
        $c_n(f) = 0$ हो वह सर्वथा शून्य है --- अर्थात् समान फूरिये गुणांकों वाले
        दो \omterm{def:b2:metric:continuity}{संतत} आवर्ती फलन मेल खाते हैं (\cref{exo:b2:fourier:4}, और अब बिना किसी स्वीकृति के)।
  \item (पारसेवाल, \omterm{def:b2:metric:continuity}{संतत} स्थिति) $S_N$ के प्रक्षेप-गुणधर्म (\cref{prop:b2:fourier:bessel}) तथा $\sigma_Nf \in \mathcal T_{N-1} \subseteq \mathcal
        T_N$ का
        उपयोग करके सिद्ध कीजिए
        \[
        \norm{f - S_Nf}_2 \leq \norm{f - \sigma_Nf}_2
        \leq \norm{f - \sigma_Nf}_\infty
        \xrightarrow[N\to\infty]{} 0 ,
        \]
        और प्रत्येक \emph{\omterm{def:b2:metric:continuity}{संतत}} $2\pi$-आवर्ती $f$ के लिए पारसेवाल की
        सर्वसमिका निष्कर्ष रूप में निकालिए।
  \item (पारसेवाल, खंडशः \omterm{def:b2:metric:continuity}{संतत} स्थिति) खंडशः \omterm{def:b2:metric:continuity}{संतत} $f$ और $\varepsilon > 0$ दिए हों तो
        $\norm{f - g}_2 \leq
        \varepsilon$ वाला कोई \omterm{def:b2:metric:continuity}{संतत} आवर्ती $g$ बनाइए \emph{(उछालों के चारों ओर
        अति छोटे अंतरालों पर $f$ के स्थान पर आनत अंतर्वेशन रखिए)}, और
        $\norm{f - S_Nf}_2 \to 0$ निष्कर्ष रूप में निकालिए \emph{(बेसेल का उपयोग कीजिए: $\norm{S_Nh}_2 \leq \norm h_2$)}:
        इस प्रकार पारसेवाल \cref{thm:b2:fourier:parseval} में कही गई पूरी व्यापकता में सत्य है ---
        और ``स्वीकृत'' हट गया।
  \item (फेयेर के लिए कोई गिब्स नहीं) \cref{exo:b2:fourier:8} से विरोध कीजिए: दिखाइए कि वर्ग
        तरंग $f$ के लिए प्रत्येक $N$ और सर्वत्र $\abs{\sigma_N(f)} \leq 1$ --- अर्थात्
        चेज़ारो औसत उस अतिक्रमण को मिटा देता है जो $S_N$ को सताता है। एक
        वाक्य में बताइए कि $F_N$ का कौन-सा गुणधर्म इसके लिए उत्तरदायी है।
\end{enumerate}

\medskip
\noindent\textbf{भाग III --- दरें।}
\begin{enumerate}[resume]
  \item $0 < \abs u \leq
        \pi$ के लिए कर्नेल के दोनों परिबंध सिद्ध कीजिए:
        \[
        F_N(u) \leq N,
        \qquad
        F_N(u) \leq \frac{\pi^2}{N u^2}
        \]
        \emph{(पहले के लिए आगमन से $\abs{\sin N\theta} \leq
        N\abs{\sin\theta}$; और दूसरे के लिए $\intcc{0}{\pi}$ पर $\sin\frac u2 \geq \frac{u}{\pi}$)}।
  \item प्रथम-आघूर्ण आकलन
        \[
        \frac{1}{2\pi}\int_{-\pi}^{\pi}\abs u\,F_N(u)\,\dd u
        \;\leq\; \frac{C\,\ln N}{N}
        \qquad (N \geq 2)
        \]
        किसी स्पष्ट अचर के साथ निष्कर्ष रूप में निकालिए \emph{($\abs u =
        \frac1N$ पर
        बाँटिए)}।
  \item निष्कर्ष निकालिए: यदि $f$ $L$-लिप्शिट्स और $2\pi$-आवर्ती हो, तो
        \[
        \norm{\sigma_N f - f}_\infty \leq \frac{C\,L\ln N}{N} .
        \]
  \item (संतृप्ति) $e_1(t) =
        \eu^{\iu t}$ के लिए $\sigma_N(e_1)$ परिकलित कीजिए और दिखाइए कि $\norm{\sigma_Ne_1 - e_1}_\infty
        = \frac1N$:
        अर्थात् सबसे चिकने फलनों के लिए भी फेयेर $\frac1N$ से तेज़ अभिसरित नहीं
        होता --- ठीक फलन-अनुक्रम अध्याय की सप्ताहांत समस्या वाली
        बर्नस्टाइन संतृप्ति का अनुरूप।
  \item (स्थानीयकरण) दिखाइए: यदि (खंडशः \omterm{def:b2:metric:continuity}{संतत}) $f$ $\intoo{x - \delta}{x + \delta}$ पर लुप्त हो, तो
        $\sigma_N(f)(x) \to 0$, चाहे $f$ अन्यत्र कितना ही उग्र क्यों न हो --- अर्थात् $x$
        पर माध्यों का अभिसरण $x$ के पास वाले $f$ को ही देखता है।
\end{enumerate}

\medskip
\noindent\textbf{भाग IV --- वाइल की समवितरण प्रमेय।} $\intco{0}{1}$ में कोई अनुक्रम
$(x_n)_{n\geq1}$ \emph{समवितरित} कहलाता है जब प्रत्येक अंतराल $\intcc{a}{b} \subseteq \intcc{0}{1}$ के लिए
\[
\frac{\#\{n \leq N : x_n \in \intcc ab\}}{N}
\xrightarrow[N\to\infty]{} b - a .
\]
\begin{enumerate}[resume]
  \item दिखाइए कि $(x_n)$ तभी समवितरित हो जाता है जब प्रत्येक \emph{\omterm{def:b2:metric:continuity}{संतत}}
        $1$-आवर्ती $f$ के लिए $\frac1N\sum_{n\leq N}f(x_n) \to \int_0^1 f$ हो \emph{($\intcc ab$ के सूचक को ऐसे दो \omterm{def:b2:metric:continuity}{संतत}
        खंडशः-आनत फलनों के बीच निचोड़िए जिनके समाकल $\varepsilon$ से भिन्न हों)}।
  \item (वाइल की कसौटी, पर्याप्तता) मान लीजिए
        \[
        \frac{1}{N}\sum_{n=1}^{N}\eu^{2\iu\pi kx_n}
        \xrightarrow[N\to\infty]{} 0
        \qquad\text{प्रत्येक } k \in \Z\setminus\{0\} .
        \text{ के लिए}\]
        दिखाइए कि पहले त्रिकोणमितीय बहुपदों के लिए $\frac1N\sum f(x_n) \to \int_0^1f$, और फिर प्रश्न 6
        (आवर्तकाल $1$ तक ले जाकर) से सभी \omterm{def:b2:metric:continuity}{संतत} $1$-आवर्ती $f$ के लिए:
        अतः प्रश्न 16 के साथ $(x_n)$ समवितरित है।
  \item मान लीजिए $\alpha$ अपरिमेय है और $x_n = \{n\alpha\}$ (भिन्नात्मक भाग)। गुणोत्तर योग
        \[
        \Bigl|\sum_{n=1}^{N}\eu^{2\iu\pi kn\alpha}\Bigr|
        \leq \frac{2}{\abs{1 - \eu^{2\iu\pi k\alpha}}}
        \qquad (k \neq 0),
        \]
        को परिबद्ध कीजिए और \emph{वाइल की प्रमेय} निष्कर्ष रूप में
        निकालिए: $(\{n\alpha\})$ $\intco{0}{1}$ में समवितरित है।
  \item इससे निष्कर्ष निकालिए कि अपरिमेय $\alpha$ के लिए $(\{n\alpha\})$ $\intcc{0}{1}$ में सघन
        है, और एक वाक्य में समझाइए कि समवितरण सघनता से यथार्थतः अधिक
        मज़बूत क्यों है।
  \item (अग्र अंक) सिद्ध कीजिए कि ऐसे पूर्णांकों $n \leq N$ का अनुपात, जिनके
        लिए $2^n$ का अग्र (दशमलव) अंक $1$ हो, $\log_{10}2 \approx 0.301$ की ओर जाता है
        \emph{(अग्र अंक $1$ का अर्थ है $\{n\log_{10}2\} \in
        \intco{0}{\log_{10}2}$; दिखाइए कि $\log_{10}2$ अपरिमेय है)}।
\end{enumerate}

\medskip
\noindent\textbf{भाग V --- समपरिमापी असमिका।} मान लीजिए $\Gamma$ लंबाई $L$ का
कोई संवृत सरल $C^1$ वक्र है जो चिह्नित क्षेत्रफल $A$ घेरता है, और उसका
प्राचलन मापित चापलंबाई से किया गया है: $z(t) =
x(t) + \iu y(t)$, $2\pi$-आवर्ती, जहाँ $\abs{z'(t)} =
\frac{L}{2\pi}$ अचर है;
और घिरा हुआ क्षेत्रफल
\[
A = \frac12\int_0^{2\pi}\bigl(x\,y' - y\,x'\bigr)\dd t
= \frac{1}{2}\,\Im\int_0^{2\pi}\conj{z}\,z'\,\dd t
\]
है (यहाँ इसे ही चिह्नित क्षेत्रफल की परिभाषा माना गया है; बहु समाकलों वाला
अध्याय ग्रीन के सूत्र से सिद्ध करता है कि वह सहज परिभाषा से मेल खाता है)।
\begin{enumerate}[resume]
  \item $z(t) = \sum_{n\in\Z}c_n\eu^{\iu nt}$ का प्रसार कीजिए ($C^1$ फलन की श्रेणी, जो \omterm{def:b2:funcseq:series}{सामान्य रूप से}
        अभिसारी है) और $z'$ पर लगाई गई पारसेवाल से सिद्ध कीजिए:
        \[
        \frac{L^2}{2\pi} = \int_0^{2\pi}\abs{z'}^2\dd t
        = 2\pi\sum_{n\in\Z}n^2\abs{c_n}^2 .
        \]
  \item इसी प्रकार $A = \pi\sum_{n\in\Z}
        n\,\abs{c_n}^2$ सिद्ध कीजिए \emph{(पारसेवाल का ध्रुवित रूप: $\frac{1}{2\pi}\int\conj f g = \sum
        \conj{c_n(f)}c_n(g)$,
        जिसे $f = z$, $g = z'$ पर लगाइए)}।
  \item (हुर्विट्स) निष्कर्ष निकालिए:
        \[
        L^2 - 4\pi A = 4\pi^2\sum_{n\in\Z}
        (n^2 - n)\abs{c_n}^2 \;\geq\; 0 ,
        \]
        और समता तभी जब $z(t) = c_0 +
        c_1\eu^{\iu t}$ --- अर्थात् कोई वृत्त। यही समपरिमापी असमिका
        है: लंबाई $L$ के संवृत वक्रों में केवल वृत्त $\frac{L^2}{4\pi}$ क्षेत्रफल
        घेरता है।
  \item तर्कसंगति की जाँचें: त्रिज्या $R$ के वृत्त के लिए समता और भुजा
        $a$ के वर्ग के लिए यथार्थ असमिका सत्यापित कीजिए; समझाइए कि
        प्रत्येक पूर्णांक $n$ के लिए (ऋणात्मक सहित) $n^2 - n \geq 0$ क्यों, और
        प्राचलन की अचर चाल कहाँ काम आई।
  \item संश्लेषण। एक-एक वाक्य में: (क) $F_N$ का वह अकेला गुणधर्म जिससे
        भाग I से III तक बहते हैं और जिसकी $D_N$ में कमी है; (ख) यहाँ का
        चेज़ारो योग घात-श्रेणी अध्याय की सप्ताहांत समस्या (फ्रोबेनियस) से
        कैसे जुड़ता है; (ग) कौन-सा लाभ केवल वाइरश्ट्रास (प्रश्न 6) से
        मिला और किसे पूरी पारसेवाल चाहिए थी; (घ) तृतीय वर्ष का खंड क्या
        जोड़ता है, एक वाक्य में ($L^2$ पूर्णता: हिल्बर्ट आधार के रूप में
        फूरिये श्रेणी)।
\end{enumerate}
\end{problem}
